\subsection{TENSOR PRODUCTS OF PRODUCTS AND DIRECT SUMS}

Let \((\mathbf{E}_{\lambda})_{\lambda \in L}\) be a family of right A-modules, \((\mathbf{F}_{\mu})_{\mu \in M}\) a family of left A-modules and consider the product modules \(C = \prod_{\lambda \in L} E_{\lambda}\) , \(D = \prod_{\mu \in M} F_{\mu}\) . The mapping \(((x_{\lambda}), (y_{\mu})) \mapsto (x_{\lambda} \otimes y_{\mu})\) of \(C \times D\) into the product Z-module

\[
\prod_ {(\lambda , \mu) \in L \times M} (E _ {\lambda} \otimes_ {A} F _ {\mu})
\]

is Z-bilinear and obviously satisfies conditions (1) (no. 1). Thus there exists (no. 1, Proposition 1) a Z-linear mapping, called canonical

\[
f: \left(\prod_ {\lambda \in L} E _ {\lambda}\right) \otimes_ {A} \left(\prod_ {\mu \in M} F _ {\mu}\right)\rightarrow \prod_ {(\lambda , \mu) \in L \times M} (E _ {\lambda} \otimes_ {A} F _ {\mu})\tag{22}
\]

such that \(f((x_{\lambda}) \otimes (y_{\mu})) = (x_{\lambda} \otimes y_{\mu})\) .

When \(C = R^{L}\) , \(D = S^{M}\) , R (resp. S) being a right (resp. left) A-module, the canonical mapping (22) associates with every tensor product \(u \otimes v\) , where u is a mapping of L into R and v a mapping of M into S, the mapping \((\lambda, \mu) \mapsto u(\lambda) \otimes v(\mu)\) of \(L \times M\) into \(R \otimes_{A} S\) ; even in this case the canonical mapping (22) is in general neither injective nor surjective (Exercise 3; cf.~Corollary 3 to Proposition 7).

When the \(E_{\lambda}\) are \(((B_{i}^{\prime}); A, (C_{j}^{\prime}))\) -multimodules and the \(F_{\mu}(A, (B_{h}^{\prime\prime}); (C_{k}^{\prime\prime}))\) -multimodules, the homomorphism (22) is also a homomorphism for the \(((B_{i}^{\prime}), (B_{h}^{\prime\prime}); (C_{j}^{\prime}), (C_{k}^{\prime\prime}))\) -multimodule structures of the two sides.

Consider now the submodule \(\mathbf{E} = \bigoplus_{\lambda \in L} \mathbf{E}_{\lambda}\) (resp. \(\bigoplus_{\mu \in M} F_{\mu}\) ) of C (resp. D); the canonical injections \(\mathbf{E} \to \mathbf{C}\) , \(\mathbf{F} \to \mathbf{D}\) define canonically a Z-linear mapping \(\mathbf{E} \otimes_{\mathbf{A}} \mathbf{F} \to \mathbf{C} \otimes_{\mathbf{A}} \mathbf{D}\) which, composed with the mapping (22), gives a Z-linear mapping \(g\) of \(\mathbf{E} \otimes_{\mathbf{A}} \mathbf{F}\) into \(\prod_{\lambda, \mu} (\mathbf{E}_{\lambda} \otimes_{\mathbf{A}} F_{\mu})\) such that

\[
g ((x _ {\lambda}) \otimes (y _ {\mu})) = (x _ {\lambda} \otimes y _ {\mu});
\]

moreover, as the families \((x_{\lambda})\) and \((y_{\mu})\) have finite support, so does \((x_{\lambda} \otimes y_{\mu})\) and hence finally g is a canonical homomorphism

\[
g: \left(\bigoplus_ {\lambda \in L} E _ {\lambda}\right) \otimes_ {A} \left(\bigoplus_ {\mu \in M} F _ {\mu}\right)\rightarrow \bigoplus_ {(\lambda , \mu) \in L \times M} (E _ {\lambda} \otimes_ {A} F _ {\mu}),\tag{23}
\]

which is a multimodule homomorphism under the same conditions as (22).

PROPOSITION 7. The canonical mapping (23) is bijective.

To prove this it suffices to define a \(\mathbf{Z}\) -linear mapping \(h\) of the direct sum \(\mathbf{G} = \bigoplus_{(\lambda, \mu) \in L \times M} (\mathbf{E}_{\lambda} \otimes_{\mathbf{A}} \mathbf{F}_{\mu})\) into \(\mathbf{E} \otimes_{\mathbf{A}} \mathbf{F}\) such that \(g \circ h\) and \(h \circ g\) are the identity mappings. But, to define a \(\mathbf{Z}\) -linear mapping of \(\mathbf{G}\) into \(\mathbf{E} \otimes_{\mathbf{A}} \mathbf{F}\) , it suffices (§ 1, no. 6, Proposition 6) to define a \(\mathbf{Z}\) -linear mapping

\[
h _ {\lambda \mu}: \mathrm{E} _ {\lambda} \otimes_ {\mathrm{A}} \mathrm{F} _ {\mu} \rightarrow \mathrm{E} \otimes_ {\mathrm{A}} \mathrm{F}
\]

for every ordered pair \((\lambda, \mu)\) , and we take \(h_{\lambda \mu} = i_{\lambda} \otimes j_{\mu}\) , where \(i_{\lambda}: \mathbf{E}_{\lambda} \to \mathbf{E}\) and \(j_{\mu}: \mathbf{F}_{\mu} \to \mathbf{F}\) are the canonical injections. Then clearly \(h \circ g\) coincides with the identity mapping for the elements of the form \(\left(\sum_{\lambda} x_{\lambda}\right) \otimes \left(\sum_{\mu} y_{\mu}\right)\) which generate the \(\mathbf{Z}\) -module \(\mathbf{E} \otimes_{\mathbf{A}} \mathbf{F}\) ; similarly \(g \circ h\) coincides with the identity mapping for the elements of the form \(\sum_{\lambda, \mu} (x_{\lambda} \otimes y_{\mu})\) which generate the \(\mathbf{Z}\) -module \(\mathbf{G}\) , since for each ordered pair \((\lambda, \mu)\) the products

\[
x _ {\lambda} \otimes y _ {\mu} \quad (x _ {\lambda} \in \mathrm{E} _ {\lambda}, y _ {\mu} \in \mathrm{F} _ {\mu})
\]

generate the Z-module \(E_{\lambda} \otimes_{A} F_{\mu}\) . Whence the proposition.

Let \(u_{\lambda}: \mathbf{E}_{\lambda} \to \mathbf{E}_{\lambda}', v_{\mu}: \mathbf{F}_{\mu} \to \mathbf{F}_{\mu}'\) be A-homomorphisms; clearly the diagram

\[
\begin{array}{c} \left(\bigoplus_ {\lambda} \mathrm{E} _ {\lambda}\right) \otimes_ {\mathrm{A}} \left(\bigoplus_ {\mu} \mathrm{F} _ {\mu}\right) \longrightarrow \bigoplus_ {\lambda , \mu} (\mathrm{E} _ {\lambda} \otimes_ {\mathrm{A}} \mathrm{F} _ {\mu}) \\ (\oplus u _ {\lambda}) \otimes (\oplus v _ {\mu}) \Bigg \downarrow \qquad \qquad \qquad \qquad \qquad \qquad \qquad \qquad \qquad \qquad \qquad \qquad \qquad \qquad \qquad \qquad \qquad \qquad \qquad \oplus (u _ {\lambda} \otimes v _ {\mu}) \\ \left(\bigoplus_ {\lambda} \mathrm{E} _ {\lambda} ^ {\prime}\right) \otimes_ {\mathrm{A}} \left(\bigoplus_ {\mu} \mathrm{F} _ {\mu} ^ {\prime}\right) \longrightarrow \bigoplus_ {\lambda , \mu} (\mathrm{E} _ {\lambda} ^ {\prime} \otimes_ {\mathrm{A}} \mathrm{F} _ {\mu} ^ {\prime}) \end{array}
\]

is commutative.

COROLLARY 1. If the left A-module F admits a basis \((b_{\mu})_{\mu \in \mathbb{M}}\) , every element of \(\mathbf{E} \otimes_{\mathbf{A}} \mathbf{F}\) can be written uniquely in the form \(\sum_{\mu} (x_{\mu} \otimes b_{\mu})\) , where \(x_{\mu} \in \mathbf{E}\) and the family \((x_{\mu})\) has finite support. The \(\mathbf{Z}\) -module \(\mathbf{E} \otimes_{\mathbf{A}} \mathbf{F}\) is isomorphic to \(\mathbf{E}^{(\mathbb{M})}\) considered as a \(\mathbf{Z}\) -module.

The basis \((b_{\mu})\) defines an isomorphism of \(\mathbf{F}\) onto \(\bigoplus_{\mu \in \mathbb{M}}\mathrm{Ab}_{\mu}\) , whence there is an isomorphism \(\mathbf{E} \otimes_{\mathbf{A}}\mathbf{F} \to \bigoplus_{\mu \in \mathbb{M}}(\mathbf{E} \otimes_{\mathbf{A}}\mathbf{Ab}_{\mu})\) by virtue of Proposition 7; as \(\xi \mapsto \xi b_{\mu}\) is an isomorphism of \(\mathbf{A}_s\) onto \(\mathbf{Ab}_{\mu}\) , \(x \mapsto x \otimes b_{\mu}\) is an isomorphism of \(\mathbf{E}\) onto \(\mathbf{E} \otimes_{\mathbf{A}} (\mathbf{Ab}_{\mu})\) by virtue of Proposition 4 of no. 4, whence the corollary. If \(\mathbf{E}\) is a \(((\mathbf{B}_i^{\prime}); \mathbf{A}, (\mathbf{C}_j^{\prime}))\) -multimodule, the canonical isomorphism \(\mathbf{E} \otimes_{\mathbf{A}} \mathbf{F} \to \mathbf{E}^{(\mathbf{M})}\) is a \(((\mathbf{B}_i^{\prime}); (\mathbf{C}_j^{\prime}))\) -multimodule isomorphism.

In particular, if also E admits a basis \((a_{\lambda})_{\lambda \in L}\) , every \(z \in E \otimes_{A} F\) may be written in one and only one way in the form \(\sum_{\lambda, \mu} (a_{\lambda} \xi_{\lambda \mu}) \otimes b_{\mu}\) , where the \(\xi_{\lambda \mu}\) belong to A (and form a family of finite support); the mapping \(z \mapsto (\xi_{\lambda \mu})_{(\lambda, \mu) \in L \times M}\) is an isomorphism of \(E \otimes_{A} F\) onto \(A^{(L \times M)}\) for the Z-module structures (and even the module structures over the centre of A). More particularly:

COROLLARY 2. If E and F are two free modules over a commutative ring C and \((a_{\lambda})\) (resp. \((b_{\mu})\) ) is a basis of the C-module E (resp. F), then \((a_{\lambda} \otimes b_{\mu})\) is a basis of the C-module E \(\otimes_{C} F\) .

By an abuse of language, the basis \((a_{\lambda} \otimes b_{\mu})\) is sometimes called the tensor product of the bases \((a_{\lambda})\) and \((b_{\mu})\) .

Remark (1). Let E be a free right A-module, F a free left A-module, \((a_{\lambda})_{\lambda \in L}\) a basis of E and \((b_{\mu})_{\mu \in M}\) a basis of F. Every element \(z \in E \otimes_{A} F\) may be written uniquely as \(\sum_{\lambda} a_{\lambda} \otimes y_{\lambda}\) , where \(y_{\lambda} \in F\) , and also uniquely as \(\sum_{\mu} x_{\mu} \otimes b_{\mu}\) , where \(x_{\mu} \in E\) . If we write \(y_{\lambda} = \sum_{\mu} \eta_{\lambda \mu} b_{\mu}\) , \(x_{\mu} = \sum_{\lambda} a_{\lambda} \xi_{\lambda \mu}\) , where the \(\xi_{\lambda \mu}\) and \(\eta_{\lambda \mu}\) belong to A, then \(\xi_{\lambda \mu} = \eta_{\lambda \mu}\) for all \((\lambda, \mu)\) , for

\[
\sum_ {\lambda} \left(a _ {\lambda} \otimes \left(\sum_ {\mu} \eta_ {\lambda \mu} b _ {\mu}\right)\right) = \sum_ {\lambda , \mu} \left((a _ {\lambda} \eta_ {\lambda \mu}) \otimes b _ {\mu}\right) = \sum_ {\mu} \left(\left(\sum_ {\lambda} a _ {\lambda} \eta_ {\lambda \mu}\right) \otimes b _ {\mu}\right).
\]

COROLLARY 3. Let \((\mathbf{E}_{\lambda})_{\lambda \in L}\) be a family of right A-modules and F a free (resp. finitely generated free) left A-module. Then the canonical mapping (22)

\[
\left(\prod_ {\lambda \in L} \mathbf {E} _ {\lambda}\right) \otimes_ {\mathbf {A}} \mathbf {F} \rightarrow \prod_ {\lambda \in L} \left(\mathbf {E} _ {\lambda} \otimes_ {\mathbf {A}} \mathbf {F}\right)
\]

is injective (resp. bijective).

If \((b_{\mu})\) is a basis of \(F\) , every element of \(\left(\prod_{\lambda \in L} E_{\lambda}\right) \otimes_{A} F\) can be written uniquely as \(z = \sum_{\mu} ((x_{\lambda}^{(\mu)}) \otimes b_{\mu})\) (Corollary 1); to say that its canonical image is zero means that, for all \(\lambda \in L\) , \(\sum (x_{\lambda}^{(\mu)} \otimes b_{\mu}) = 0\) , hence \(x_{\lambda}^{(\mu)} = 0\) for all \(\lambda \in L\) and all \(\mu\) (Corollary 1) and therefore \(z = 0\) .

Showing that the canonical mapping is bijective when F admits a finite basis is immediately reduced, by virtue of Proposition 7, to the case where

\(F = A_{s}\) ; but then the two sides are canonically identified with \(\prod_{\lambda \in L} E_{\lambda}\) (no. 4, Proposition 4) and after these identifications the canonical mapping (22) becomes the identity.

COROLLARY 4. Let A be a ring with no divisor of zero, E a free right A-module and F a free left A-module. Then the relation \(x \otimes y = 0\) in \(\mathbf{E} \otimes_{\mathbf{A}} \mathbf{F}\) implies \(x = 0\) or \(y = 0\) .

Let \((a_{\lambda})\) be a basis of E, \((b_{\mu})\) a basis of F and let \(x = \sum_{\lambda} a_{\lambda} \xi_{\lambda}, y = \sum_{\mu} \eta_{\mu} b_{\mu}\) ; then \(x \otimes y = \sum_{\lambda, \mu} ((a_{\lambda} \xi_{\lambda} \eta_{\mu}) \otimes b_{\mu})\) and the relation \(x \otimes y = 0\) implies \(\xi_{\lambda} \eta_{\mu} = 0\) for every ordered pairs of indices \((\lambda, \mu)\) (Corollary 1). Hence, if \(x \neq 0\) , that is \(\xi_{\lambda} \neq 0\) for at least one \(\lambda\) , it follows that \(\eta_{\mu} = 0\) for all \(\mu\) , whence \(y = 0\) .

COROLLARY 5. Let E be a right A-module, F a left A-module, M a submodule of E and N a submodule of F. If M is a direct factor of E and N a direct factor of F, the canonical homomorphism \(M \otimes_{A} N \to E \otimes_{A} F\) is injective and the image of \(M \otimes_{A} N\) under this homomorphism is a direct factor of the Z-module \(E \otimes_{A} F\) .

This follows immediately from Proposition 7.

Note that if E is a \(((B_{i}^{\prime}); A, (C_{j}^{\prime}))\) -multimodule and F an \((A, (B_{h}^{\prime\prime}); (C_{k}^{\prime\prime}))\) -multimodule and M and N direct factors in these multimodules, \(M \otimes N\) is a direct factor of the \(((B_{i}^{\prime}), (B_{h}^{\prime\prime}); (C_{j}^{\prime}), (C_{k}^{\prime\prime}))\) -multimodule \(E \otimes F\) .

COROLLARY 6. Let P be a projective left A-module and E, F two right A-modules. For every injective homomorphism \(u: \mathbf{E} \to \mathbf{F}\) , the homomorphism

\[
u \otimes \mathrm{l} _ {\mathbf {P}}: \mathrm{E} \otimes_ {\mathbf {A}} \mathrm{P} \rightarrow \mathrm{F} \otimes_ {\mathbf {A}} \mathrm{P}
\]

is injective.

There exists a left A-module \(\mathbf{Q}\) such that \(\mathbf{L} = \mathbf{P} \oplus \mathbf{Q}\) is free (§2, no. 2, Proposition 4) and \(u \otimes 1_{\mathbf{L}}\) is identified (Proposition 7) with

\[
(u \otimes 1 _ {\mathbf {P}}) \oplus (u \otimes 1 _ {\mathbf {Q}});
\]

it therefore suffices to prove the corollary when P is free (§ 1, no. 6, Corollary 1 to Proposition 7). The same argument reduces the problem to the case where \(P = A_{s}\) , which follows immediately from no. 4, Proposition 4.

COROLLARY 7. Let C be a commutative ring. If E and F are two projective C-modules, E ⊗\_C F is a projective C-module.

This follows immediately from Corollary 5 and the fact that the tensor product of two free C-modules is a free C-module (Corollary 2).

Remark 2. Under the hypotheses of Proposition 7, let \(\mathbf{E}_{\lambda}^{\prime}\) be a submodule of \(\mathbf{E}_{\lambda}, \mathbf{F}_{\mu}^{\prime}\) a submodule of \(\mathbf{F}_{\mu}\) and let \(\mathbf{E}^{\prime} = \bigoplus_{\lambda \in \mathbf{L}} \mathbf{E}_{\lambda}^{\prime}, \mathbf{F}^{\prime} = \bigoplus_{\mu \in \mathbf{M}} \mathbf{F}_{\mu}^{\prime}\) . Let \(\operatorname{Im}(\mathbf{E}^{\prime} \otimes_{\mathbf{A}} \mathbf{F}^{\prime})\)

(resp. \(\operatorname{Im}(\mathbf{E}_{\lambda}^{\prime}\otimes_{\mathbf{A}}\mathbf{F}_{\mu}^{\prime})\) denote the image of \(\mathbf{E}^{\prime}\otimes_{\mathbf{A}}\mathbf{F}^{\prime}\) (resp. \(\mathbf{E}_{\lambda}^{\prime}\otimes_{\mathbf{A}}\mathbf{F}_{\mu}^{\prime}\) ) in \(\mathbf{E}\otimes_{\mathbf{A}}\mathbf{F}\) (resp. \(\mathbf{E}_{\lambda}\otimes_{\mathbf{A}}\mathbf{F}_{\mu}\) ) under the canonical mapping; then the isomorphism (23) identifies the sub- \(\mathbf{Z}\) -modules

\[
\operatorname{Im} \left(\mathrm{E} ^ {\prime} \otimes_ {\mathbf {A}} \mathrm{F} ^ {\prime}\right) \quad \text { and } \quad \bigoplus_ {(\lambda , \mu) \in \mathbf {L} \times \mathbf {M}} \operatorname{Im} \left(\mathrm{E} _ {\lambda} ^ {\prime} \otimes_ {\mathbf {A}} \mathrm{F} _ {\mu} ^ {\prime}\right);
\]

this follows immediately from the commutativity of the diagram

\[
\begin{array}{c} \left(\bigoplus_ {\lambda} \mathrm{E} _ {\lambda} ^ {\prime}\right) \otimes_ {\mathrm{A}} \left(\bigoplus_ {\mu} \mathrm{F} _ {\mu} ^ {\prime}\right) \longrightarrow \left(\bigoplus_ {\lambda} \mathrm{E} _ {\lambda}\right) \otimes_ {\mathrm{A}} \left(\bigoplus_ {\mu} \mathrm{F} _ {\mu}\right) \\ \Biggl \downarrow \\ \bigoplus_ {\lambda , \mu} (\mathrm{E} _ {\lambda} ^ {\prime} \otimes_ {\mathrm{A}} \mathrm{F} _ {\mu} ^ {\prime}) \longrightarrow \bigoplus_ {\lambda , \mu} (\mathrm{E} _ {\lambda} \otimes_ {\mathrm{A}} \mathrm{F} _ {\mu}) \end{array}
\]

where the vertical arrows are the canonical isomorphisms.

